%! AP Statistics — Unit 3, Lesson 3.3 — Activity KEY
\documentclass[10pt]{article}
\usepackage{apstats-article}
\usepackage{apstats-key}

\begin{document}

\pageheader{Unit 3, Lesson 3.3}{Group Activity KEY: Design Your Sample}
\namepartnerperiod

\vspace{0.1in}

\begin{practicebox}
\small
\textbf{A. SRS}
\begin{enumerate}[leftmargin=1.5em, itemsep=4pt]
  \item \ans{Assign each student a unique number 1--8{,}000.}
  \item \ans{Random number table, calculator randInt, or random number generator.}
  \item \ans{Every group of 200 students has the same chance; no systematic exclusion.}
  \item \ans{Could under-represent smaller grade levels if sizes differ substantially.}
\end{enumerate}
\end{practicebox}

\vspace{0.08in}

\begin{practicebox}
\small
\textbf{B. Stratified}
\begin{enumerate}[leftmargin=1.5em, itemsep=4pt]
  \item \ans{Four grade levels (9, 10, 11, 12).}
  \item \ans{50 from each grade (assuming equal grade sizes); adjust proportionally if sizes differ.}
  \item \ans{Ensures each grade is fairly represented; useful if satisfaction varies across grade levels.}
\end{enumerate}
\end{practicebox}

\vspace{0.08in}

\begin{practicebox}
\small
\textbf{C. Cluster}
\begin{enumerate}[leftmargin=1.5em, itemsep=4pt]
  \item \ans{Approx. 4 schools $\times$ 500 students/school = 2{,}000 students.}
  \item \ans{Yes — each school contains students from all grade levels, making it a heterogeneous cluster.}
  \item \ans{Logistically easy — survey whole schools rather than tracking down individuals.}
\end{enumerate}
\end{practicebox}

\vspace{0.08in}

\begin{practicebox}
\small
\textbf{D. Systematic}
\begin{enumerate}[leftmargin=1.5em, itemsep=4pt]
  \item \ans{$k = 8{,}000 / 200 = 40$}
  \item \ans{Randomly pick a number between 1--40; select that student, then every 40th student after.}
  \item \ans{If the list is organized in a pattern (e.g., alphabetical by grade), some grades may be over-represented.}
\end{enumerate}
\end{practicebox}

\vspace{0.08in}

\begin{practicebox}
\small
\textbf{Part 2:}
\begin{enumerate}[leftmargin=1.5em, itemsep=4pt]
  \item \ans{Stratified — ensures each grade level is proportionally represented.}
  \item \ans{Systematic or cluster — easier logistically.}
  \item \ans{Stratified random sample by grade level best balances statistical accuracy (proportional representation) with the district's interest in comparing across grades. Cluster sampling is practical but may miss grade-level nuance.}
\end{enumerate}
\end{practicebox}

\end{document}
